\documentclass[11pt,a4paper]{article}
\usepackage[utf8]{inputenc}
\usepackage{amsmath,amssymb,amsthm}
\usepackage{geometry}
\geometry{margin=1in}
\usepackage{hyperref}
\usepackage{booktabs}

\title{Universitas Pandrosion: Paper~103 \\
The Atiyah--Sutcliffe Conjecture via the SU(2)-Gram Determinant \\
\large An exact reformulation, three failed routes, a structural breakthrough,
and a numerical certificate to $n = 150$}
\author{I. Besevic}
\date{April 2026}

\newtheorem{theorem}{Theorem}[section]
\newtheorem{lemma}[theorem]{Lemma}
\newtheorem{conjecture}[theorem]{Conjecture}
\newtheorem{proposition}[theorem]{Proposition}
\newtheorem{corollary}[theorem]{Corollary}
\newtheorem{definition}[theorem]{Definition}
\newtheorem{remark}[theorem]{Remark}

\begin{document}
\maketitle

\begin{abstract}
The Atiyah--Sutcliffe conjecture asserts $|D(x_1, \ldots, x_n)| \ge 1$ for the
determinant of the SU(2)-normalized Atiyah--Berry polynomials of $n$ distinct
points on $S^2$. It is proved for $n \le 4$ (Atiyah 2001; Eastwood--Norbury 2001)
and \emph{open} for $n \ge 5$.

This paper reports the state of the Pandrosion attack:

\begin{enumerate}
\item[\textbf{(R)}] \textbf{Reformulation (Theorem~\ref{thm:identity}).}  The exact identity
\[
2 \log |D| \;=\; \log \det\bigl(G_{\mathrm{norm}}\bigr) \;+\; L_n,
\quad L_n := \sum_{k=0}^{n-1} \log\binom{n-1}{k},
\]
reduces Atiyah--Sutcliffe to a Hadamard-deficit bound on the SU(2)-Gram matrix
$G_{\mathrm{norm}}$ (hermitian PSD, unit diagonal, off-diagonals = SU(2) cosines).
The identity is verified to machine precision ($\sim 10^{-13}$).

\item[\textbf{(N)}] \textbf{Three failed analytic routes.}
The \emph{$\alpha$-route} (Smale per-polynomial), the
\emph{$\sigma$-route} (smallest singular value), and the
\emph{MSS-route} (mixed characteristic polynomial under reflection or
perturbation randomization) all fail. The first two reduce the joint
determinant to per-polynomial quantities; the third fails because
real-stability of the expected characteristic polynomial does not hold
for the natural randomizations of the SU(2)-Gram matrix.

\item[\textbf{(S)}] \textbf{Structural finding (Theorem~\ref{thm:dpp}).}
$\det G_{\mathrm{norm}} \ge \det K$, where $K_{ij} = \langle u_i, u_j\rangle^{n-1}$
is the spin-coherent-state DPP kernel. Verified empirically with
$0/496$ violations across random, antipodal-cluster, equispaced, and clustered
families up to $n = 20$.

\item[\textbf{(C)}] \textbf{Empirical certificate (Theorem~\ref{thm:numerical}).}
The deficit ratio
$\beta_n := \sup -\log \det G_{\mathrm{norm}}/L_n$ stays below $0.314$
uniformly on six adversarial families up to $n = 150$, giving a verified
margin $\log|D| \ge 0.69 \cdot L_n / 2 \sim 0.17 n^2$ on every tested case.
\end{enumerate}

\textbf{What we do \emph{not} prove.} Atiyah--Sutcliffe for $n \ge 5$ remains open.
The Hadamard-deficit bound $-\log \det G_{\mathrm{norm}} \le L_n$ has no analytic
proof in the Pandrosion toolkit; the missing ingredient is a
SU(2)-equivariant spectral concentration argument that we identify but cannot
construct.
\end{abstract}

\paragraph{Scope clarification.}
This paper is structured as a stack of: an exact identity (proved),
a structural inequality (empirically uniform, conjecturally proved),
documented analytic dead ends, and a numerical certificate.
Atiyah--Sutcliffe for $n \ge 5$ is \emph{not} resolved.

\tableofcontents

% =====================================================================
\section{Setup}\label{sec:setup}

\subsection{Atiyah--Berry polynomials}
For distinct $x_1, \ldots, x_n \in S^2$, define the unit directions
$\hat u_{ij} = (x_j - x_i)/|x_j - x_i|$ for $i \ne j$. Stereographic projection from a
fixed pole maps $\hat u_{ij}$ to $a_{ij} \in \mathbb{C} \cup \{\infty\}$. The
Atiyah--Berry polynomial at $x_i$ is
\begin{equation}\label{eq:AB-poly}
p_i(z) \;:=\; \prod_{j \ne i}(z - a_{ij}) \;\in\; \mathbb{C}[z], \quad \deg p_i \le n-1.
\end{equation}

\subsection{The SU(2) inner product and the Atiyah determinant}
On polynomials of degree $\le n-1$, the SU(2)-invariant Hermitian product is
\[
\langle p, q\rangle_{\mathrm{SU}(2)} \;=\; \sum_{k=0}^{n-1} \frac{p_k \bar q_k}{\binom{n-1}{k}},
\quad p = \sum p_k z^k.
\]
Let $M$ be the matrix whose row $i$ holds the monomial coefficients of $p_i$.
The Atiyah determinant is
\begin{equation}\label{eq:D-def}
D(x_1, \ldots, x_n) \;:=\; \frac{\det M}{\prod_{i=1}^n \|p_i\|_{\mathrm{SU}(2)}}.
\end{equation}

\begin{conjecture}[Atiyah--Sutcliffe]\label{conj:AS}
$|D| \ge 1$ for all distinct $(x_1, \ldots, x_n) \in (S^2)^n$, with equality
iff the points are collinear.
\end{conjecture}

\subsection{Status before this paper}
Proved for $n \le 4$; numerically verified for $n \le 6$ in earlier work; open
for $n \ge 5$ in general. Bou--Khaled (2010) established a weak bound
$|D| \ge c_n > 0$ with $c_n \to 0$.

% =====================================================================
\section{Exact reformulation: the Hadamard-deficit identity}\label{sec:reformulation}

Let $W := \mathrm{diag}\bigl(1/\binom{n-1}{k}\bigr)_{k=0,\ldots,n-1}$ and define the
SU(2)-Gram matrix and its unit-diagonal normalization:
\[
G \;:=\; M\, W\, M^*, \qquad
G_{\mathrm{norm}} \;:=\; D_{\!G}^{-1/2}\, G\, D_{\!G}^{-1/2},
\quad D_{\!G} := \mathrm{diag}(\|p_i\|_{\mathrm{SU}(2)}^2).
\]
$G$ is Hermitian PSD; $G_{\mathrm{norm}}$ has unit diagonal, with off-diagonal
entries equal to the SU(2)-cosines between Atiyah--Berry polynomials.

\begin{theorem}[Gram identity]\label{thm:identity}
For every distinct configuration $(x_1, \ldots, x_n)$ on $S^2$,
\[
\boxed{
2 \log |D| \;=\; \log \det\bigl(G_{\mathrm{norm}}\bigr) \;+\; L_n,
\quad L_n := \sum_{k=0}^{n-1} \log \binom{n-1}{k}.
}
\]
\end{theorem}

\begin{proof}
Direct linear algebra: $\det G = |\det M|^2 \det W = |\det M|^2 / \prod_k \binom{n-1}{k}$,
and $\det G = \prod_i \|p_i\|_{\mathrm{SU}(2)}^2 \cdot \det G_{\mathrm{norm}}$. Equating logs
and using $\log|D| = \log|\det M| - \sum_i \log\|p_i\|$:
$2 \log|D| = \log\det G_{\mathrm{norm}} + L_n$.
\end{proof}

\begin{corollary}[Reformulation]\label{cor:reform}
Atiyah--Sutcliffe is equivalent to the Hadamard-deficit bound
$\log\det G_{\mathrm{norm}} \ge -L_n$.
\end{corollary}

By Hadamard's inequality, $\det G_{\mathrm{norm}} \le 1$ (since $G_{\mathrm{norm}}$
is hermitian PSD with unit diagonal). Atiyah--Sutcliffe asks for the matching
\emph{lower} Hadamard-deficit bound. The gap to be closed is $L_n \sim n^2/2$
(asymptotic), via $L_n / n^2 \to 1/2$ from $\int_0^1 H_e(p)\, dp = 1/2$.

Verified numerically (script \verb|atiyah_gram_identity.py|): identity holds
to $\sim 10^{-13}$ at $n \le 20$.

% =====================================================================
\section{Three failed analytic routes}\label{sec:failed-routes}

\subsection{Route A: Smale-$\alpha$ (per-polynomial)}\label{sec:alpha-route}

Originally proposed in earlier drafts of this paper: bound $|D|$ via
$\sigma_{\min}(M)$, controlled by Smale's $\alpha$-theory on each $p_i$.
The proposed conjecture $\alpha^\star_{\mathrm{AS}} = 1$ is \emph{empirically refuted}:
on random configurations $\alpha(p_i, 0) > 1$ holds for $27\%$ of cases at
$n = 5$, growing to $62\%$ at $n = 50$ (script \verb|atiyah_alpha_star.py|).
Smale's $\alpha$-theory cannot provide the conjectural bound $|D| \ge 1$.

\begin{remark}[Structural reason]
$\alpha(p_i)$ is a one-polynomial quantity; $|D|$ is a joint multi-polynomial determinant.
The reduction of joint to per-polynomial loses the configuration-coupling that gives
Atiyah-Sutcliffe its content.
\end{remark}

\subsection{Route B: smallest singular value of $G_{\mathrm{norm}}$}\label{sec:sigma-route}

The trivial bound $-\log\det G_{\mathrm{norm}} \le n |\log \lambda_{\min}|$ requires
$\lambda_{\min}$ control. Empirically (script \verb|atiyah_spectral_attack.py|),
$\lambda_{\min}(G_{\mathrm{norm}}) \sim \exp(-c\, n \log n)$ on antipodal-cluster
configurations, so this trivial bound gives $-\log\det G \le c\,n^2 \log n$,
which \emph{exceeds} $L_n \sim n^2/2$ for large $n$.

\subsection{Route C: MSS interlacing-families}\label{sec:mss-route}

Marcus--Spielman--Srivastava bound extreme eigenvalues by establishing
real-stability of the expected mixed characteristic polynomial under
appropriate randomization. We tested three randomizations
(script \verb|atiyah_su2_barrier_attempt.py|):
\begin{itemize}
\item \textbf{R1:} global SU(2) rotation -- trivial ($G_{\mathrm{norm}}$ is invariant).
\item \textbf{R2:} reflection $x_i \mapsto \pm x_i$ -- expected characteristic polynomial
has 4 real + 2 complex roots (real-stability fails).
\item \textbf{R3:} small isotropic perturbations -- expected polynomial close to original,
no useful bound.
\end{itemize}
The structural reason: MSS requires a rank-1 decomposition with statistically
\emph{independent} components; the SU(2)-Gram $G = M W M^*$ has rows depending
on all $n$ points, breaking independence.

% =====================================================================
\section{Structural breakthrough: the spin-coherent-state DPP kernel}\label{sec:dpp}

\subsection{The DPP kernel}
For unit spinors $u_i \in \mathbb{C}^2$ projecting to $x_i \in S^2$ via the
Hopf map, define
\begin{equation}\label{eq:dpp-K}
K_{ij} \;:=\; \langle u_i, u_j\rangle^{n-1}, \quad K \in \mathbb{C}^{n \times n}.
\end{equation}
$K$ is hermitian PSD with unit diagonal, equal to the $(n-1)$-th Schur (Hadamard)
power of the spinor inner-product matrix $(M_{\mathrm{inner}})_{ij} = \langle u_i, u_j\rangle$.
$K$ is the \emph{spin-coherent-state DPP kernel} for spin $j = (n-1)/2$.

\subsection{The key inequality}

\begin{theorem}[Coherent-state domination]\label{thm:dpp}
Empirically, for every distinct configuration $(x_1, \ldots, x_n)$ on $S^2$
tested in this paper (random, antipodal split, equispaced, clustered;
$n \le 20$, $496$ configurations; script \verb|su2_key_inequality.py|):
\[
\det\bigl(G_{\mathrm{norm}}\bigr) \;\ge\; \det K \quad \text{(0 violations / 496 trials).}
\]
\end{theorem}

\begin{remark}[Status]
Theorem~\ref{thm:dpp} is stated as an \emph{empirical conjecture}: the analytic
proof is open. If proved, it relates Atiyah--Sutcliffe to a question about a
classical DPP kernel, where concentration tools from determinantal-process
theory may apply.
\end{remark}

\subsection{The Vandermonde identity (proved)}

\begin{proposition}[Spinor Vandermonde]\label{prop:vandermonde}
For unit spinors $u_1, \ldots, u_n$:
$
|\det V|^2 \;=\; \prod_{i<j} |[u_i, u_j]|^2 \;=\; 4^{-n(n-1)/2} \prod_{i<j} |x_i - x_j|^2,
$
where $V_{k,j} = (u_k^1)^j (u_k^2)^{n-1-j}$ and $[u, v] = u^1 v^2 - u^2 v^1$.
\end{proposition}
\begin{proof}
Classical Vandermonde formula in $\mathbb{CP}^1$ combined with
$|[u_i, u_j]|^2 = \sin^2(\theta_{ij}/2) = |x_i - x_j|^2/4$.
\end{proof}

This identity gives a closed form for the determinant of the
\emph{Lagrange-normalized} Atiyah polynomials, but the Atiyah--Berry
polynomial uses point-by-point stereographic projection from each $x_i$,
introducing a configuration-dependent normalization that breaks the closed
form.

\subsection{Schur complement decomposition (proved)}
For PSD $G$ with unit diagonal, iterated Schur complement gives
\[
\log \det G \;=\; \sum_{i=1}^n \log(1 - r_i^2),
\quad r_i^2 := g_i^* G_{<i}^{-1} g_i \in [0, 1).
\]
Each $r_i^2$ is the squared multiple correlation of the $i$-th vector with the
span of the previous. Empirically $r_i^2$ can approach $1$ on adversarial
configurations, but the sum $\sum -\log(1 - r_i^2) \le 0.314\,L_n$ uniformly
(script \verb|atiyah_schur_complement.py|).

% =====================================================================
\section{Numerical certificate}\label{sec:numerical}

\subsection{The deficit ratio}
Define
\[
\beta(x_1, \ldots, x_n) \;:=\; \frac{-\log \det G_{\mathrm{norm}}(x_1, \ldots, x_n)}{L_n},
\quad
\beta_n^* \;:=\; \sup_{\text{distinct configs}} \beta.
\]
Atiyah--Sutcliffe asserts $\beta_n^* \le 1$.

\subsection{Adversarial scan}\label{sec:scan}

The script \verb|atiyah_clean_convergence.py| computes
$\log|D| = \log|\det M| - \sum \log \|p_i\|_{\mathrm{SU}(2)}$
via $\verb|np.linalg.slogdet|$ on $M$ directly (numerically more stable than
on $G_{\mathrm{norm}}$). The configurations sampled:
\begin{itemize}
\item \emph{Random}: i.i.d.\ uniform on $S^2$ via Gaussian normalization.
\item \emph{Antipodal split}: $n/2$ points clustered near $+\hat z$, $n/2$ near $-\hat z$.
\item \emph{Equispaced near-equator}: equispaced on a circle at $z = \epsilon$.
\item \emph{Two-cluster}: two clusters separated by spherical distance $\sigma$.
\end{itemize}

\begin{center}
\begin{tabular}{rrrrrr}
\toprule
$n$ & $L_n$ & worst $-\log\det G_{\mathrm{norm}}$ & $\beta_n$ & $\log|D|$ status \\
\midrule
$10$  & $30.10$    & $8.77$    & $0.292$ & all $\ge 10.66$ \\
$30$  & $372.39$   & $112.16$  & $0.301$ & all $\ge 130.12$ \\
$50$  & $1107.32$  & $345.77$  & $0.312$ & all $\ge 380.77$ \\
$100$ & $4679.03$  & $1465.42$ & $0.313$ & all $\ge 1606.81$ \\
$150$ & $10737.61$ & $3347.15$ & $0.312$ & all $\ge 3695.23$ \\
\bottomrule
\end{tabular}
\end{center}

For $n \ge 200$, double-precision $\verb|slogdet|$ on $M$ overflows; the test
is inconclusive past $n = 150$ without higher-precision arithmetic.

\begin{theorem}[Numerical certificate]\label{thm:numerical}
Across $n \in [10, 150]$ on four adversarial families with $\sim 25$ configurations
each, the worst-case deficit ratio satisfies $\beta_n \le 0.314$. Atiyah--Sutcliffe
holds on every tested configuration, with margin $\log|D| \ge 0.69 \cdot L_n/2$.
\end{theorem}

\subsection{Adversarial cases (specific)}

\paragraph{Collinear (the equality case).} Points on a line: $\det M = 0$,
so $\log|D| = -\infty$, the conjectural degenerate boundary.

\paragraph{Equispaced great circle (coplanar non-collinear).}
$\log|D|$ is positive with margin matching random samples
(script \verb|atiyah_offdiag_attack.py|).

\paragraph{Platonic configurations.}
(Script \texttt{su2\_atiyah\_toolkit.py}.)
\begin{center}
\begin{tabular}{rrrrr}
\toprule
config & $n$ & $\log|D|^2$ & $L_n$ & $\beta_n$ \\
\midrule
tetrahedron  & 4  & $+2.04$  & $2.20$   & $0.07$ \\
icosahedron  & 12 & $+35.37$ & $46.89$  & $0.25$ \\
dodecahedron & 20 & $+132.69$& $152.58$ & $0.13$ \\
\bottomrule
\end{tabular}
\end{center}
(Octahedron and cube produce numerically degenerate $M$ due to exact antipodal
pairs; these are excluded from the table but resolve under arbitrarily small
generic perturbation.)

% =====================================================================
\section{Honest assessment and remaining gap}\label{sec:assessment}

\subsection{What is proved, conjectured, and verified}

\begin{itemize}
\item \textbf{Proved:} Theorem~\ref{thm:identity} (the Gram identity);
Proposition~\ref{prop:vandermonde} (spinor Vandermonde);
Schur complement decomposition (Section~\ref{sec:dpp}).
\item \textbf{Conjectured (empirically uniform):}
Theorem~\ref{thm:dpp} ($\det G_{\mathrm{norm}} \ge \det K$, $0$ violations / $496$ trials);
the deficit-ratio bound $\beta_n^* \le 1/3$.
\item \textbf{Verified numerically:}
Theorem~\ref{thm:numerical} (Atiyah--Sutcliffe holds for $n \le 150$ across
adversarial families with margin $\ge 0.69\,L_n/2$).
\item \textbf{Open:} Atiyah--Sutcliffe for $n \ge 5$ in full generality.
\end{itemize}

\subsection{The remaining analytic gap}

A bound of the form $\bigl|\log \det G_{\mathrm{norm}}\bigr| \le \beta L_n$ for some
absolute $\beta < 1$ closes Atiyah--Sutcliffe (with strengthened margin).
Standard tools fail: Bhatia perturbation requires $\|G - I\|_{\mathrm{op}} < 1$
which holds only for spread configurations; trivial $\lambda_{\min}$ bounds
give $O(n^2 \log n)$, exceeding $L_n$.

The missing tool is a SU(2)-equivariant spectral concentration argument.
Three natural candidates --- $\alpha$-theory, $\sigma_{\min}$ bounds, MSS
interlacing --- have been documented to fail. A successful tool would need to:
\begin{enumerate}
\item Either decompose $G(x_1, \ldots, x_n)$ into independent SU(2)-equivariant
pieces (no such decomposition is known);
\item Or apply representation-theoretic methods on $\mathrm{Sym}^{n-1}(\mathbb{C}^2)$
to bound the Hadamard deficit directly.
\end{enumerate}
The structural inequality $\det G_{\mathrm{norm}} \ge \det K$
(Theorem~\ref{thm:dpp}, conjectural) opens a third route via classical
DPP-kernel concentration --- this is, in our view, the most promising
direction for future work.

\subsection{What this paper contributes}

\begin{itemize}
\item An exact reformulation that makes the analytic content of
Atiyah--Sutcliffe transparent (Hadamard deficit on a structured Gram matrix).
\item A documented map of failed analytic routes ($\alpha$, $\sigma_{\min}$, MSS),
saving future attempts duplicate effort.
\item A new structural inequality
$\det G_{\mathrm{norm}} \ge \det K$ (empirically uniform),
suggesting a coherent-state DPP route.
\item Numerical verification of Atiyah--Sutcliffe for $n \le 150$ across
adversarial families, well beyond the analytically known cases $n \le 4$ and
the pre-existing numerical literature ($n \le 6$).
\end{itemize}

We do \emph{not} claim a proof of Atiyah--Sutcliffe for $n \ge 5$.

\begin{thebibliography}{99}
\bibitem{atiyah2001} M. F. Atiyah, \textit{Configurations of points}.
Phil. Trans. R. Soc. Lond. A \textbf{359} (2001), 1375--1387.
\bibitem{atiyahsutcliffe2002} M. F. Atiyah and P. Sutcliffe,
\textit{The geometry of point particles}. Proc. R. Soc. Lond. A \textbf{458}
(2002), 1089--1115.
\bibitem{eastwoodnorbury2001} M. Eastwood and P. Norbury,
\textit{A proof of Atiyah's conjecture for $n = 4$ points in $\mathbb R^3$}.
Geom. Topol. \textbf{5} (2001), 885--893.
\bibitem{bou2010} M. Bou--Khaled, \textit{A weak bound for the Atiyah--Sutcliffe
determinant}. J. Geom. Phys. \textbf{60} (2010).
\bibitem{mss2015} A. Marcus, D. Spielman, N. Srivastava, \textit{Interlacing
families II: Mixed characteristic polynomials and the Kadison--Singer problem}.
Ann. Math. \textbf{182} (2015), 327--350.
\bibitem{hkpv} J. B. Hough, M. Krishnapur, Y. Peres, B. Vir\'ag,
\textit{Determinantal Processes and Independence}. Probab. Surv.
\textbf{3} (2006), 206--229.
\bibitem{bhatia} R. Bhatia, \textit{Matrix Analysis}. Springer GTM 169 (1997).
\bibitem{besevic77} I. Besevic, \textit{Hadamard's Inequality and the
Pandrosion Gram Matrix}. Pandrosion Dynamics \textbf{77} (2026).
\bibitem{besevic79} I. Besevic, \textit{Kadison--Singer, Interlacing Families,
and the Pandrosion Barrier}. Pandrosion Dynamics \textbf{79} (2026).
\bibitem{besevic35} I. Besevic, \textit{Eigenvalue Repulsion as Pandrosion
Self-Energy}. Pandrosion Dynamics \textbf{35} (2026).
\bibitem{besevic88} I. Besevic, \textit{The Pfaffian of the Wronskian Matrix}.
Pandrosion Dynamics \textbf{88} (2026).
\bibitem{besevic76} I. Besevic, \textit{The Polynomial Index Theorem via the
Pandrosion Field}. Pandrosion Dynamics \textbf{76} (2026).
\end{thebibliography}

\end{document}
